% =============================================================================
% CAPÍTULO 4: Marco Metodológico
% =============================================================================
\chapterUnnumbered{Capítulo 4: Marco Metodológico}

[En un párrafo explica brevemente los contenidos de este capítulo antes
de la distribución en distintos apartados.

Describe el método que has seguido para desarrollar el proyecto. Este
apartado debe contener las principales características de la
investigación desarrollada, en cualquier tipo de TFM. Se debe describir
el proceso que implica la detección de la necesidad o interés del
trabajo, la búsqueda de la información y las fuentes documentales
empleadas, el tipo de investigación (cualitativa, cuantitativa o
sociocrítica, de tipo descriptiva, observación, analítica, estudio de
casos, intervención, etnográfica, comparada, proyecto,
investigación-acción, buena práctica, etc.). Se trata de etiquetar y
catalogar el proceso, diseño y tipo de investigación que has
desarrollado.

En este punto, explicaremos cómo se ha llevado a cabo la investigación
del Trabajo Fin de Máster, los pasos que hemos seguido, los materiales,
recursos que hemos diseñado o utilizado para proceder a la recogida de
la información y la puesta en práctica de todo lo que se haya diseñado
para tal fin.

Los apartados a considerar (aunque puedes variarlo para adecuarlo a tu
proyecto concreto, ampliando o reduciendo apartados) son los
siguientes:]

\setcounter{section}{0}
\section{Población y muestra}

[Describe la población objeto de estudio y el muestreo realizado para
obtener la muestra que representa dicha población.]

\section{Objetivos de investigación}

[¿Qué pretendes evaluar? Si los objetivos son los mismos que los
objetivos del TFM no hace falta que incluyas este apartado.]

\section{Variables intervinientes e instrumentos de evaluación}

[Describe qué vas a medir y a través de qué instrumentos]

\section{Diseño experimental}

[Describe el método de investigación elegido y su procedimiento]

\section{Modelo de análisis de datos}

[Describe cómo vas a analizar los datos que obtengas en tu
investigación. Si los vas a analizar cuantitativa o cualitativamente. Si
vas a realizar análisis estadísticos y mediante qué programas
informáticos, etc.].

[Se recomienda un tamaño de 5 a 15 páginas para este capítulo]
